\[
x(t)=u(t)-2u(t-2)+u(t-5)
\]
\[
h(t)=e^{2t}u(1-t)
\]

ابتدا \(x(t)\) را به صورت قطعه‌ای می‌نویسیم:
\[
x(t)=
\begin{cases}
0 & t<0\\
1 & 0\le t<2\\
-1 & 2\le t<5\\
0 & t\ge 5
\end{cases}
\]

کانولوشن برابر است با:
\[
y(t)=x(t)*h(t)=\int_{-\infty}^{\infty}x(\tau)h(t-\tau)\,d\tau
\]

اکنون:
\[
h(t-\tau)=e^{2(t-\tau)}u(1-t+\tau)
\]

از شرط تابع پله داریم:
\[
1-t+\tau \ge 0
\quad \Longrightarrow \quad
\tau \ge t-1
\]

پس:
\[
y(t)
=e^{2t}\int_{t-1}^{\infty}x(\tau)e^{-2\tau}\,d\tau
\]

با توجه به بازه‌های غیرصفر \(x(\tau)\)، چهار حالت به دست می‌آید.

برای \(t<1\)، کل بازه‌ی غیرصفر \(x(\tau)\) داخل انتگرال قرار می‌گیرد:
\[
y(t)
=e^{2t}
\left(
\int_{0}^{2}e^{-2\tau}\,d\tau
-
\int_{2}^{5}e^{-2\tau}\,d\tau
\right)
\]
\[
y(t)
=
\frac{e^{2t}}{2}
\left(
1-2e^{-4}+e^{-10}
\right)
\]

برای \(1\le t<3\)، کران پایین انتگرال داخل بازه‌ی \(0\le \tau<2\) قرار می‌گیرد:
\[
y(t)
=e^{2t}
\left(
\int_{t-1}^{2}e^{-2\tau}\,d\tau
-
\int_{2}^{5}e^{-2\tau}\,d\tau
\right)
\]
\[
y(t)
=
\frac{1}{2}
\left(
e^2-2e^{2t-4}+e^{2t-10}
\right)
\]

برای \(3\le t<6\)، فقط بخشی از بازه‌ی منفی \(x(\tau)\) باقی می‌ماند:
\[
y(t)
=
-e^{2t}
\int_{t-1}^{5}e^{-2\tau}\,d\tau
\]
\[
y(t)
=
\frac{1}{2}
\left(
e^{2t-10}-e^2
\right)
\]

برای \(t\ge 6\)، دیگر هیچ هم‌پوشانی‌ای بین سیگنال‌ها وجود ندارد:
\[
y(t)=0
\]

در نتیجه:
\[
\boxed{
y(t)=
\begin{cases}
\dfrac{e^{2t}}{2}\left(1-2e^{-4}+e^{-10}\right) & t<1\\[6pt]
\dfrac{1}{2}\left(e^2-2e^{2t-4}+e^{2t-10}\right) & 1\le t<3\\[6pt]
\dfrac{1}{2}\left(e^{2t-10}-e^2\right) & 3\le t<6\\[6pt]
0 & t\ge 6
\end{cases}
}
\]
